\chapter{Descripción del Trabajo}
\label{cap:descripcionTrabajo}

El presente trabajo consiste en el desarrollo de una plataforma web dirigida a nutricionistas, con el objetivo de facilitar la gestión de sus clientes, la planificación de dietas personalizadas y el seguimiento continuo de los tratamientos nutricionales.

\section{Objetivo del Proyecto}

La plataforma tiene como finalidad centralizar todas las herramientas necesarias para que un nutricionista pueda gestionar de forma profesional y eficiente su trabajo diario, mejorando la comunicación con los clientes y automatizando tareas repetitivas.

\section{Funcionalidades de la Plataforma}

A continuación se detallan las principales funcionalidades desarrolladas:

\subsection{Funcionalidades para Nutricionistas}

\begin{itemize}
  \item Gestión de perfiles de clientes incluyendo historial clínico, objetivos, medidas y hábitos.
  \item Clasificación de clientes por categorías personalizadas de objetivos (pérdida de peso, ganancia muscular, mantenimiento) y por tipo de dieta (cetogénica, alta en proteínas, baja en carbohidratos, etc.).
  \item Filtros de búsqueda para segmentación rápida de clientes según las categorías establecidas.
  \item Constructor de dietas personalizadas con plantillas reutilizables y cálculo automático de macronutrientes.
  \item Biblioteca de alimentos organizada mediante etiquetas (vegano, sin gluten, etc.).
  \item Zona privada de preparación de dietas.
  \item Historial completo de dietas previas asignadas.
  \item Exportación de dietas e informes en formato PDF.
  \item Inteligencia Artificial personalizada para asistir en la generación de dietas según parámetros definidos por el profesional.
  \item Análisis predictivo del progreso estimando el peso o la evolución futura de cada cliente en base al plan asignado.
  \item Creación de cuestionarios personalizados para recopilar información de hábitos, actividad física y otros datos relevantes.
  \item Gestión de la agenda de citas, disponibilidad y confirmación online o presencial.
  \item Recordatorios automáticos por correo, SMS o notificaciones.
  \item Programación de mensajes de seguimiento posteriores a las citas.
  \item Chat privado con cada cliente y posibilidad de adjuntar archivos.
  \item Zona de recursos profesionales (vídeos, guías, plantillas, eBooks) para compartir y vender entre nutricionistas.
  \item Gestión de grupos de clientes para retos o programas colectivos.
  \item Control de roles de usuarios: nutricionistas, administrativos y otros colaboradores.
  \item Panel de métricas detalladas del progreso de cada cliente y estadísticas globales de la actividad profesional.
\end{itemize}

\subsection{Funcionalidades para Clientes}

\begin{itemize}
  \item Acceso completo a sus dietas, datos personales y evolución registrada.
  \item Registro de peso, medidas, fotos de progreso y analíticas.
  \item Envío de feedback sobre bienestar, adherencia, energía y dificultades.
  \item Cuestionarios de satisfacción periódicos con notificación automática al nutricionista en caso de incidencias.
  \item Pantalla resumen diaria (modo "Hoy") mostrando la dieta del día, próximas citas y mensajes pendientes.
  \item Solicitud de citas directamente desde el calendario de disponibilidad del profesional.
  \item Selección de modalidad de cita (presencial u online).
  \item Recepción de recordatorios, avisos y mensajes automatizados o personalizados.
  \item Comunicación directa con el nutricionista mediante chat.
\end{itemize}

\subsection{Funcionalidades Generales}

\begin{itemize}
  \item Plataforma web completamente responsive, compatible tanto en escritorio como en dispositivos móviles.
  \item Sistema seguro de gestión de roles y permisos de acceso.
  \item Automatización de procesos rutinarios (mensajería, notificaciones, recordatorios, etc.).
  \item Exportación automática de documentación en formato PDF.
  \item Arquitectura escalable apta tanto para nutricionistas independientes como para clínicas de mayor tamaño.
\end{itemize}
